\input{../tex-inputs/latex-leseansicht-vorspann}

\section[Arthur und Heinrich Schnitzler sowie Paul Marx an Gustav Schwarzkopf, 15. 3. 1925]{L04362 Arthur und Heinrich Schnitzler sowie Paul Marx an Gustav Schwarzkopf, 15. 3. 1925}
\nopagebreak\mylabel{L04362v}
\rehead{ }\normalsize\beginnumbering\briefempfaengerindex{Schwarzkopf, Gustav@\textsc{Schwarzkopf, Gustav}!zzzMarx, Paul@\emph{von Paul Marx}!1925-03-151@{15. 3. 1925}|(be}\briefempfaengerindex{Schwarzkopf, Gustav@\textsc{Schwarzkopf, Gustav}!zzzSchnitzler, Heinrich@\emph{von Heinrich Schnitzler}!1925-03-151@{15. 3. 1925}|(be}\briefempfaengerindex{Schwarzkopf, Gustav@\textsc{Schwarzkopf, Gustav}!zzzSchnitzler, Arthur@\emph{von Arthur Schnitzler}!1925-03-151@{15. 3. 1925}|(be}
\toendnotes[C]{\smallbreak\pagebreak[2]}
\correspDesc{Versand  durch Arthur Schnitzler, Heinrich Schnitzler, Paul Marx am 15. 3. 1925 in Berlin
\newline{}Übermittlung  am 16. 3. 1925 in Berlin
\newline{}Erhalt  durch Gustav Schwarzkopf im Zeitraum [16. 3. 1925
                  – 20. 3. 1925?] in Wien}\toendnotes[C]{\smallbreak}
\Standort{DLA, A:Schnitzler, HS.1985.1.1897.}
\physDesc{Bildpostkarte, 391 Zeichen
\newline{}Handschrift Arthur Schnitzler: Bleistift, lateinische Kurrent
\newline{}Handschrift Heinrich Schnitzler: Bleistift, lateinische Kurrent
\newline{}Handschrift Paul Marx: Bleistift, deutsche Kurrent
\newline{}Versand: Stempel: »\nobreak{}\oindex{Berlin@\textbf{Berlin}, \emph{Hauptstadt}|pwk}Berlin W 9, 16. 3. 25, 7–8V\nobreak{}«.  }\toendnotes[C]{\smallbreak}\pstart{}{\pb}Herrn Guſtav Schwarzkopf\pend{}\pstart{}Wien I\oindex{I., Innere Stadt@\textbf{I., Innere Stadt}, \emph{Verwaltungsgebiet}|pw}\pend{}\pstart{}Tiefer Graben 17\oindex{Wien@\textbf{Wien}!I., Innere Stadt@\textbf{I., Innere Stadt}!Tiefer Graben 17@\textbf{Tiefer Graben 17}, \emph{Wohngebäude}|pw}\pend{}{\bigskip}
\pstart
           \noindent{}\centering{}{\pb}\textcolor{gray}{\textbf{Hotel Esplanade, Berlin}}\oindex{Hotel Esplanade [Berlin]@\textbf{Hotel Esplanade [Berlin]}, \emph{Hotel}|pw}\pend
           \vspace{1em}
\pstart
           \raggedleft{}{\pb}Berlin\oindex{Berlin@\textbf{Berlin}, \emph{Hauptstadt}|pw}{ }15. 3. 25 \pend
           \vspace{0.5em}
\pstart
           Am 14. war hier Première\eventindex{Lessing-Theater@\textbf{Lessing-Theater}!Premiere von Das Märchen, 14.3.1925@Premiere von Das Märchen, 14.3.1925|pwv} des Märchen\pwindex{Schnitzler, Arthur 15. 5. 1862 Wien – 21. 10. 1931 ebd.@\textsc{Schnitzler, Arthur} (15. 5. 1862 Wien – 21. 10. 1931 ebd.), \emph{Schriftsteller, Mediziner}!Märchen. Schauspiel in drei Aufzügen@\strich\emph{Das Märchen. Schauspiel in drei Aufzügen}|pw} (am Lessingtheater\orgindex{Lessing-Theater@Lessing-Theater|pw}) – 32 Jahre nach Wien\oindex{Wien@\textbf{Wien}, \emph{Verwaltungsgebiet}|pw}\eventindex{Volkstheater@\textbf{Volkstheater}!Uraufführung von Das Märchen, 1.12.1893@Uraufführung von Das Märchen, 1.12.1893|pwv}. Ich sah mir am gleichen Abend Charleys Tante\pwindex{Thomas, Brandon 25.\,12.\,1856 Liverpool – 19.\,6.\,1914@\textsc{Thomas, Brandon} (25.\,12.\,1856 Liverpool – 19.\,6.\,1914), \emph{Schriftsteller}!Charley’s Aunt. A Farce in Three Acts@\strich\emph{Charley’s Aunt. A Farce in Three Acts}|pw}\eventindex{Schauspielhaus Berlin@\textbf{Schauspielhaus Berlin}!Aufführung von Charley’s Tante, 14.3.1925@Aufführung von Charley’s Tante, 14.3.1925|pwv} an \strikeout{\textcolor{gray}{mit}} mit Heini als Jack\pwindex{Thomas, Brandon 25.\,12.\,1856 Liverpool – 19.\,6.\,1914@\textsc{Thomas, Brandon} (25.\,12.\,1856 Liverpool – 19.\,6.\,1914), \emph{Schriftsteller}!Charley’s Aunt. A Farce in Three Acts@\strich\emph{Charley’s Aunt. A Farce in Three Acts}|pwv} (singend und tanzend.) Wer mir das am
                  \label{K_L04362-1v}\edtext{1. Dezember 1893}{\lemma{\textnormal{\emph{1. Dezember 1893}}}\Cendnote{\textnormal{Dem Tag der Uraufführung\eventindex{Volkstheater@\textbf{Volkstheater}!Uraufführung von Das Märchen, 1.12.1893@Uraufführung von Das Märchen, 1.12.1893|pwk} am \emph{Deutschen
                     Volkstheater}XXXX ORGangabe fehlt in Wien\oindex{Wien@\textbf{Wien}, \emph{Verwaltungsgebiet}|pwk}.}}}\label{K_L04362-1}
               prophezeiht hätte. Heute trat  Heini\pwindex{Schnitzler, Heinrich 9.\,8.\,1902 Hinterbrühl – 12.\,7.\,1982 Wien@\textsc{Schnitzler, Heinrich} (9.\,8.\,1902 Hinterbrühl – 12.\,7.\,1982 Wien), \emph{Regisseur, Schauspieler}|pw} zum 500. Mal auf! –\pend
           \pstart Herzliche Grüße. Ihr \spacefill\mbox{Arthur}\pend{}\selectlanguage{ngerman}\vspace{1em}
\pstart
           \noindent{}{[}hs. Schnitzler:{]} Viele herzliche Grüße\pend
           \pstart \spacefill\mbox{Heini.}\pend{}\selectlanguage{ngerman}\vspace{1em}
\pstart
           \noindent{}{[}hs. Marx:{]} Mit beſten Empfehlungen Ihr\pend
           \pstart \spacefill\mbox{PaulMarx}\pend{}\selectlanguage{ngerman}\endnumbering\briefempfaengerindex{Schwarzkopf, Gustav@\textsc{Schwarzkopf, Gustav}!zzzMarx, Paul@\emph{von Paul Marx}!1925-03-151@{15. 3. 1925}|)be}\briefempfaengerindex{Schwarzkopf, Gustav@\textsc{Schwarzkopf, Gustav}!zzzSchnitzler, Heinrich@\emph{von Heinrich Schnitzler}!1925-03-151@{15. 3. 1925}|)be}\briefempfaengerindex{Schwarzkopf, Gustav@\textsc{Schwarzkopf, Gustav}!zzzSchnitzler, Arthur@\emph{von Arthur Schnitzler}!1925-03-151@{15. 3. 1925}|)be}\mylabel{L04362h}
\begin{anhang}
\end{anhang}\newcommand{\dateiname}{L04362}\newcommand{\titel}{Arthur und Heinrich Schnitzler sowie Paul Marx an Gustav Schwarzkopf, 15. 3. 1925}\newcommand{\editorInnen}{Herausgegeben von Jahnke, SelmaMüller, Martin Anton}\input{../tex-inputs/latex-leseansicht-abspann}
